\documentclass{beamer}
\usepackage[T1]{fontenc}
\usepackage[dvipsnames]{xcolor}
\usepackage{graphicx, float}
\graphicspath{{res/}}
\usepackage{amsfonts,amsmath}
\usepackage{listings}
\usepackage{pifont}
\usepackage{hyperref}

\usetheme{moderno}

\lstset{
  basicstyle=\ttfamily\footnotesize,
  keywordstyle=\color{blue},
  commentstyle=\color{gray},
  stringstyle=\color{red},
  numbers=left,
  numberstyle=\tiny\color{gray},
  breaklines=true,
  frame=single
}

\newcommand{\dingcheck}{{\color{moderngreen}\ding{52}}}
\newcommand{\dingcross}{{\color{modernred}\ding{54}}}

\title{Debug Information Validation for Optimized Code}
\subtitle{PLDI '20 -- Yuanbo Li, Shuo Ding, Qirun Zhang, Davide Italiano}
\author{Silas Kraume}
\institute{Heinrich-Heine-Universit\"at D\"usseldorf}
\date{08.07.2026}

\begin{document}
\maketitle


\begin{frame}[noframenumbering]{Agenda}
  \tableofcontents
\end{frame}

\section{Motivation}


\begin{frame}{Warum optimierten Code debuggen?}
  \begin{itemize}
    \item Fast alle Produktionssoftware wird mit Optimierungen kompiliert
    \item Industriebeispiele:
      \begin{itemize}
        \item GNU Autotools: Standard \texttt{-g -O2} in Release-Builds
        \item Red Hat Developer Guide: Empfiehlt \texttt{-g -O2} f\"ur alle Entwickler
        \item \dots
      \end{itemize}
    \item Post-Mortem Debugging: Analyse von Coredumps nach Crashes in Produktion
    \item Performance-Bugs: Treten nur in optimiertem Code auf
    \item Historischer Kontext: Vor 20 Jahren: ``Optimierungen ausschalten zum Debuggen''  Heute: Debuggen von optimiertem Code ist obligatorisch
  \end{itemize}
\end{frame}
% Dies ist nicht theoretisch—große Software wird mit Debug-Info + Optimierungen ausgeliefert.
% Bei Crashes in Produktion müssen Entwickler optimierte Binärdateien debuggen.
% Falsche Debug-Info macht Diagnose unmöglich oder irreführend.
% Diese Arbeit schließt eine kritische Lücke in der Software-Qualitätssicherung.

\begin{frame}{Was fehlt im aktuellen Testing?}
  \begin{itemize}
    \item Bisheriger Testing-Fokus:
      \begin{itemize}
        \item[\dingcheck] Compiler-Korrektheit (Maschinencode-Ausgabe)
        \item[\dingcheck] Debugger-Funktionalit\"at (interaktives Testing)
        \item[\dingcross] Korrektheit der Debug-Informationen f\"ur optimierten Code
      \end{itemize}
    \item Das Problem: Compiler k\"onnen generieren:
      \begin{itemize}
        \item[\dingcheck] Korrekten Maschinencode
        \item[\dingcross] Falsche Debug-Informationen
      \end{itemize}
    \item Konsequenz: Debugger zeigt falsche Werte, obwohl Programm korrekt l\"auft
    \item Forschungsfrage: Generieren Compiler immer korrekte Debug-Informationen f\"ur optimierten Code?
  \end{itemize}
\end{frame}
% Dies ist die zentrale Erkenntnis—Maschinencode und Debug-Info können divergieren.
% Traditionelles Compiler-Testing erkennt dies nicht.
% Das Programm läuft korrekt, aber der Debugger lügt.
% Dies ist extrem gefährlich für Post-Mortem-Debugging.
% Dieses Paper schließt diese Lücke mit dem ersten systematischen Framework.

\section{Hintergrund: Debug-Informationen}


\begin{frame}{Was sind Debug-Informationen?}
  \begin{itemize}
    \item Zweck: Stellt Verbindung zwischen Quellcode und Maschinencode her
    \item Generiert von: Compilern mit \texttt{-g} Flag
    \item Formate:
      \begin{itemize}
        \item DWARF (Linux/Unix -- am h\"aufigsten, Fokus dieses Papers)
        \item PDB (Windows)
        \item STABS (\"Altere Unix-Systeme)
      \end{itemize}
    \item Verwendet von: Debuggern (GDB, LLDB) zur Interpretation laufender Programme
    \item Enth\"alt: Funktionen, Variablen, Zeilennummern, Typen, Speicherorte
    \item Kritisch: Debugger verlassen sich vollst\"andig auf diese Informationen zur Anzeige von Variablenwerten
  \end{itemize}
\end{frame}
% Debug-Informationen sind die Brücke zwischen Quellcode und Executable.
% DWARF ist das Standardformat für Unix-ähnliche Systeme.
% Ohne sie können Debugger keine Variablennamen oder Quellzeilen anzeigen.
% Der Compiler generiert dies zusammen mit Maschinencode.
% Bei Optimierung muss diese Zuordnung korrekt aktualisiert werden.

\begin{frame}[fragile]{DWARF Debug-Information Beispiel}
  \begin{lstlisting}[basicstyle=\ttfamily\tiny]
DW_TAG_compile_unit // Kompilierungseinheit-Metadaten (``Wurzelknoten'')
  DW_AT_producer ("clang ver 7.0.1")
  DW_AT_language (DW_LANG_C99)
  DW_AT_name ("a.c")
  DW_AT_low_pc (0x00400480) // Startadresse der Funktion
  DW_AT_high_pc (<offset-from-lowpc>50) // Endadresse der Funktion

DW_TAG_subprogram // Funktionsinformationen
  DW_AT_name ("main")
  DW_AT_decl_line (2)
  DW_AT_frame_base (DW_OP_reg6 RBP)

DW_TAG_variable // Variableninformationen
  DW_AT_name ("a")
  DW_AT_location (DW_OP_fbreg -8) // Offset von RBP (Variablenort)
  DW_AT_decl_line (1) // Deklarationszeile im Quellcode
  DW_AT_type (0x00000060 "int") // Typ-Referenz (z.B. int)

DW_TAG_variable // Variableninformationen
  DW_AT_name ("b")
  DW_AT_location (DW_OP_fbreg -12)
  DW_AT_decl_line (1)
  DW_AT_type (0x00000060 "int")
  \end{lstlisting}
\end{frame}
% DWARF ist ein Baum von Einträgen, jeder mit Tags und Attributen.
% Variablenorte werden mit Debug-Operationen (DW_OP_*) spezifiziert.
% Bei Optimierung müssen diese Orte korrekt aktualisiert werden.
% Dies ist Abbildung 2 aus dem Paper—tatsächliche DWARF-Ausgabe von Clang.
% DW_OP_fbreg -8 bedeutet: Variable 'a' ist bei Frame-Base-Register-Offset -8.
% Bei Optimierung müssen diese Orte aktualisiert werden.
% Bugs entstehen, wenn der Ort auf falschen Speicher oder veraltete Werte zeigt.
% Der Debugger liest dies, um Variablenwerte anzuzeigen.

\section{Problemstellung}


\begin{frame}{Problemstellung -- Kernherausforderung}
  \begin{itemize}
    \item Kernproblem:
      \begin{itemize}
        \item Generieren von geeigneten Eingabeprogrammen $P$ f\"ur systematisches Debug-Info-Testing
      \end{itemize}
    \vspace{0.5cm}
    \item Warum existierende Techniken versagen:
      \begin{itemize}
        \item Compiler-Testing generiert Programme mit festen Ausgabe-Statements (print, return)
        \item Programme werden nur einmal ausgef\"uhrt bzgl. Programmeingabe
        \item Debugging ist interaktiver -- Benutzer inspizieren mehrere Programmpunkte mit vielen Variablen
      \end{itemize}
  \end{itemize}
\end{frame}
% Existierende Compiler-Testing-Techniken sind unzureichend für Debug-Info-Validierung.
% Der interaktive Charakter des Debuggens erfordert andere Testgenerierung.
% Compiler-Optimierungen machen dies besonders herausfordernd.
% Dies motiviert die Notwendigkeit von Actionable Programs.

\begin{frame}{Herausforderungen bei Debug-Info-Validierung}
  \begin{itemize}
    \item 1. Code Location Problem:
      \begin{itemize}
        \item Compiler-Optimierung rearrangiert Instruktionen und eliminiert Statements
        \item Viele Programmpunkte im Quellcode k\"onnten wegoptimiert werden
        \item Breakpoints an diesen Orten stoppen Programmausf\"uhrung nicht notwendigerweise
      \end{itemize}
    \vspace{0.5cm}
    \item 2. Data Value Problem:
      \begin{itemize}
        \item Debugger k\"onnen jede Variable an einem spezifischen Programmpunkt inspizieren
        \item Ausgabe einer beliebigen Variable kann undefined behaviors (UB) verursachen
        \item Beispiele: Zugriff auf uninitialisierte Variable, Out-of-Bound-Array-Element
      \end{itemize}
  \end{itemize}
\end{frame}

\section{Formale Definitionen}


\begin{frame}{Definition -- Actionable Program}
  \begin{block}{Definition (Actionable Program)}
    F\"ur einen spezifischen Optimierungsprozess enth\"alt ein \textbf{actionable program}
    $P_{\langle s,v\rangle}$ einen Programmort $s$ und eine Variable $v$, sodass sowohl
    im unoptimierten Programm als auch im optimierten Programm $P'_{\langle s,v\rangle}$ der Debugger an $s$ stoppen
    kann und die Ausgabe von $v$ an $s$ wohldefiniert ist (kein undefined behaviour).
  \end{block}
  \vspace{0.2cm}
  \begin{itemize}
    \item Menge der Actionable Programs:
      \[
        |\mathcal{P}| = \sum_{i \in I} s_i \cdot v_i
      \]
      Wobei:
      \begin{itemize}
        \item $I$ = Menge der ausgef\"uhrten Zeilen
        \item $s_i$ = Programmort
        \item $v_i$ = Anzahl der wohldefinierten Variablen an jedem Ort
      \end{itemize}
  \end{itemize}
\end{frame}


\begin{frame}{Definition -- Korrektheit}
  \begin{itemize}
    \item Value-Printing Sequence: $A_{\langle s,v\rangle}$ repr\"asentiert \\
      \texttt{break(s)} $\rightarrow$ \texttt{run} $\rightarrow$ \texttt{print(v)} $\rightarrow$ \texttt{quit}
    \item Debugger-Ausgabe: $\mathit{DBG}(P, A_{\langle s,v\rangle})$ gibt den Wert von $v$ nach Abschluss der Debugging-Session aus
    \item Referenzwert: $\mathit{PRINT}(P, s, v)$ bezeichnet den ausgegebenen Referenzwert
      \begin{itemize}
        \item Print-Statement einf\"ugen, um Wert von $v$ an $s$ im Quellcode anzuzeigen
        \item Programm $P$ ohne Optimierung kompilieren und ausf\"uhren
      \end{itemize}
  \end{itemize}
  \vspace{0.3cm}
  \begin{block}{Definition (Korrektheit)}
    Die Debug-Informationen f\"ur optimierten Code $P'$ sind korrekt gdw.
    \[
      \mathit{DBG}(P', A_{\langle s,v\rangle}) \equiv \mathit{PRINT}(P, s, v)
    \]
    wobei $P$ ein unoptimiertes Programm und $P'$ ein optimiertes Programm aus demselben Quellcode ist.
  \end{block}
\end{frame}

\section{Konkrete Bug-Beispiele}

\begin{frame}{Definition -- Korrektheit anwenden -- Konkrete Szenarien}
  \begin{center}
    \footnotesize
    \begin{tabular}{|l|c|c|c|}
      \hline
      \textbf{Szenario} & $\mathit{PRINT}(P,s,v)$ & $\mathit{DBG}(P', A_{\langle s,v\rangle})$ & \textbf{Urteil} \\
      \hline
      Korrekte Debug-Info & 5 & 5 & \dingcheck\ Korrekt \\
      \hline
      Falsche Debug-Info & 5 & 1 & \dingcross\ BUG ($5 \neq 1$) \\
      \hline
      Variable Optimized Out & 5 & null & \dingcheck\ Akzeptabel \\
      \hline
      Debugger Crash & 5 & crash & \dingcross\ BUG (Debugger sollte nicht crashen) \\
      \hline
    \end{tabular}
  \end{center}
  \vspace{0.3cm}
  \begin{itemize}
    \item Wichtig: ``Optimized out'' ist kein Bug
    \item Bug-Kriterium: Nur falsche Werte (nicht null) verletzen Korrektheit
    \item Paper-Fokus: Testing der Korrektheit, nicht Vollst\"andigkeit
  \end{itemize}
\end{frame}


\begin{frame}[fragile]{LLVM Bug 43893 (Clang-9)}
  \begin{columns}
    \column{0.55\textwidth}
      \begin{lstlisting}[firstnumber=3,basicstyle=\ttfamily\small]
int c = --a;
int l = -8L;
l = c;       // optimized out
c > (b = 0); // l
      \end{lstlisting}
      Erwartet: Variable \texttt{l} sollte \texttt{-1} in Zeile 6 sein \\
      Tats\"achlich (LLDB, -g -O2): Zeigt \texttt{-8}

    \column{0.42\textwidth}
      DWARF Bug: InstCombine-Pass verwirft ungenutzte Variable, \\
      markiert Debug-Info aber nicht als unavailable \\
      \vspace{0.2cm}
      Korrektheitsverletzung: \\
      $\mathit{PRINT}(P,\text{Zeile }6,l) = -1$ \\
      $\mathit{DBG}(P', A_{\langle 6,l\rangle}) = -8$ \\
      $-1 \not\equiv -8 \rightarrow$ \alert{BUG} (Definition verletzt) \\
      \vspace{0.2cm}
      Status: In Development-Version gefixt
  \end{columns}
\end{frame}
% Der Maschinencode ist korrekt (printf würde -1 zeigen).
% Aber LLDB zeigt -8, weil DWARF-Info nicht aktualisiert wurde.
% Der InstCombine-Optimierungs-Pass entfernte die Variable, ließ aber veraltete Debug-Info.
% Dies ist genau der Bug-Typ, den dieses Framework erkennt.

\begin{frame}[fragile]{GCC Bug 89463 (GCC-8.3)}
  \begin{columns}
    \column{0.55\textwidth}
      \begin{lstlisting}[firstnumber=3,basicstyle=\ttfamily\small]
int i;
for (; a < 10; a++)
    i = 0;
for (; i < 6; i++)
    ;
opt_me_not(); // i
      \end{lstlisting}
      Erwartet: Variable \texttt{i} sollte \texttt{6} in Zeile 8 sein \\
      Tats\"achlich (GDB, -g -O3): Zeigt \texttt{0} (falsch!)

    \column{0.42\textwidth}
      Ursache: Dead Code Elimination bereinigt \\
      degenerierte Phi-Nodes nicht \\
      \vspace{0.2cm}
      Korrektheitsverletzung: \\
      $\mathit{PRINT}(P,\text{Zeile }8,i) = 6$ \\
      $\mathit{DBG}(P', A_{\langle 8,i\rangle}) = 0$ \\
      $6 \not\equiv 0 \rightarrow$ \alert{BUG} (Definition verletzt) \\
      \vspace{0.2cm}
      Status: In Development-Version gefixt
  \end{columns}
\end{frame}
% Die Schleife wird optimiert, aber Debug-Info sagt i = 0.
% Phi-Nodes in SSA-Repräsentation wurden nicht korrekt bereinigt.
% Dies verhinderte Aktualisierung der Debug-Informationen.
% Wiederum: Maschinencode korrekt, nur Debug-Info falsch.

\section{Methodik: Shadow Validation \& Line Augmentation}

\begin{frame}[fragile]{Actionable Program Transformation}

  \begin{itemize}
    \item Originalprogramm $P$:
    \begin{lstlisting}[]
char a,b;
int main() {
    unsigned c;
    --b;
    c=2;
    if (a)
        b=0;
    return 0;
}
  \end{lstlisting}
    \item Actionable Program $P_{\langle 6,c\rangle}$: Breakpoint in Zeile 6, inspiziere Variable \texttt{c} \dingcheck
    \item Actionable Program $P_{\langle 5,b\rangle}$: F\"uge \texttt{opt\_me\_not()} in Zeile 5 ein, inspiziere \texttt{b} \dingcheck
    \item Non-Actionable Program $P_{\langle 4,c\rangle}$: Variable \texttt{c} uninitialisiert in Zeile 4 \dingcross
  \end{itemize}
\end{frame}
% Das Framework bestimmt automatisch, welche Paare gültig sind.

\begin{frame}{Systematisches Testing Framework \"Ubersicht}
  \textbf{Eingabe:} Wohlgeformtes Seed-Programm $P_w$
  \vspace{0.3cm}
  \begin{itemize}
    \item Schritt 1: Pre-Processing
      \begin{itemize}
        \item Sammle Line Coverage Information $S$
        \item Baue Variablenmengen $V_s$ f\"ur jede Zeile (AST-Traversierung)
      \end{itemize}
    \item Schritt 2: Actionable Program Generierung
      \begin{itemize}
        \item Shadow Validation (l\"ose Data Value Problem)
        \item Line Augmentation (l\"ose Code Location Problem)
        \item Generiere Menge der actionable programs $\mathcal{P}$
      \end{itemize}
    \item Schritt 3: Systematisches Testing
      \begin{itemize}
        \item Kompiliere unoptimiert $\rightarrow$ Erhalte Referenz $\mathit{PRINT}(P,s,v)$
        \item Kompiliere optimiert $\rightarrow$ Erhalte Debugger-Ausgabe $\mathit{DBG}(P', A_{\langle s,v\rangle})$
        \item Vergleiche Ausgabe $\rightarrow$ Reporte Bug wenn $\mathit{DBG} \neq \mathit{PRINT}$ und $\mathit{DBG} \neq \texttt{null}$
      \end{itemize}
  \end{itemize}
\end{frame}


\begin{frame}{Shadow Validation Technik (Data Value Problem)}
  \begin{itemize}
    \item Ziel: Sicherstellen, dass Inspizieren von Variable $v$ an Zeile $s$ sicher ist (kein UB)
    \item Methode:
      \begin{itemize}
        \item Erstelle ``Shadow Program'' $P_1$ durch Einf\"ugen von \texttt{print(v)} an Zeile $s$
        \item F\"uhre $P_1$ mit dynamischer Analyse aus
        \item Erkenne undefined behaviors (uninitialisierte Reads, Out-of-Bounds, etc.)
      \end{itemize}
    \item Ergebnis:
      \begin{itemize}
        \item Wenn UB erkannt $\rightarrow$ Kandidat verwerfen (non-actionable)
        \item Wenn sicher $\rightarrow$ Fahre fort mit Line Augmentation
      \end{itemize}
    \item Tools: CompCert Interpreter + Clang UndefinedBehaviorSanitizer
  \end{itemize}
\end{frame}
% CompCert ist ein verifizierter Compiler—sein Interpreter ist zuverlässig für UB-Erkennung (beweist formal die korrektheit (mathematisch, formale methodik)).
% Clang's Sanitizer bietet Runtime-Monitoring.
% Dies sichert, dass alle Testfälle wohlgeformt sind.

\begin{frame}[fragile]{Shadow Validation Beispiel}
  \begin{itemize}
        \begin{lstlisting}[firstnumber=3,basicstyle=\ttfamily\tiny]
  unsigned c;
  --b;
  c=2; // c hier initialisiert
        \end{lstlisting}
    \item Shadow Program $P_1$ (Invalide):
        \begin{lstlisting}[firstnumber=3,basicstyle=\ttfamily\tiny]
  unsigned c;
  --b;
  printf("\%d", c); // UB! c hier uninitialisiert
  c=2; // c hier initialisiert
        \end{lstlisting}
        \begin{itemize}
          \item Abgelehnt durch Shadow Validation (non-actionable)
        \end{itemize}
    \item Shadow Program $P_1$ (Valide):
        \begin{lstlisting}[firstnumber=3,basicstyle=\ttfamily\tiny]
  unsigned c;
  --b;
  c=2; // c hier initialisiert
  printf("\%d", c); // sicher, c hier initialisiert
        \end{lstlisting}
        \begin{itemize}
          \item Akzeptiert durch Shadow Validation (actionable)
        \end{itemize}
  \end{itemize}
\end{frame}


\begin{frame}{Line Augmentation Technik (Code Location Problem)}
  \begin{itemize}
    \item Ziel: Sicherstellen, dass Debugger an Zeile $s$ stoppen kann, auch wenn optimiert
    \item Methode:
      \begin{itemize}
        \item Pr\"ufe DWARF Line Number Table f\"ur Zeile $s$
        \item Wenn Zeile $s$ fehlt/optimiert:
          \begin{itemize}
            \item F\"uge Aufruf von unoptimierbarer Funktion \texttt{opt\_me\_not()} vor Zeile $s$ ein
            \item Dies wirkt als Barriere f\"ur Optimierung
            \item F\"uge Meta-Information (Kommentare) hinzu, um Variable $v$ zu spezifizieren
          \end{itemize}
      \end{itemize}
    \item Implementierung: \texttt{opt\_me\_not()} in separater Kompilierungseinheit (verhindert Inlining)
    \item Alternative: \texttt{\_\_asm\_\_ volatile("" ::: "memory")} als zus\"atzliche Barriere
  \end{itemize}
\end{frame}
% Die Funktion ist extern, Compiler kann sie nicht wegoptimieren.
% Wir können auch Inline-Assembly als zusätzliche Barriere verwenden.
% Dies zwingt Compiler, den Ort für Debugging zu behalten.
% Ohne dies wären viele Bugs nicht erkennbar.
% Dies ermöglicht Testing von Zeilen, die sonst unerreichbar wären.

\section{Testalgorithmus}


\begin{frame}[fragile]{Systematischer Testalgorithmus}
  \begin{block}{Systematischer Testalgorithmus}
    \textbf{Eingabe:} Ein wohlgeformtes Testprogramm $P_w$ \\
    \begin{enumerate}
      \item $\mathcal{P} \leftarrow$ \textsc{ActionableProgramGeneration}($P_w$)
      \item \textbf{for each} $P_{\langle s,v\rangle} \in \mathcal{P}$ \textbf{do}
        \begin{enumerate}
          \item $A_{\langle s,v\rangle} \leftarrow$ \textsc{ExtractActions}($P_{\langle s,v\rangle}$)
          \item $P' \leftarrow$ kompiliere $P_{\langle s,v\rangle}$ mit Optimierung und ``\texttt{-g}''
          \item \textbf{if} $\mathit{DBG}(P', A_{\langle s,v\rangle})$ ist nicht null \textbf{and} \\
            $\mathit{DBG}(P', A_{\langle s,v\rangle}) \neq \mathit{PRINT}(P, s, v)$ \textbf{then}
            \begin{enumerate}
              \item Reporte Bug mit Programm $P_{\langle s,v\rangle}$, \\
                Action $A_{\langle s,v\rangle}$ und verwendetem Optimierungs-Flag
            \end{enumerate}
        \end{enumerate}
    \end{enumerate}
  \end{block}
\end{frame}

\section{Evaluation \& Ergebnisse}


\begin{frame}{Bug-Zusammenfassung}
  \begin{center}
    \begin{tabular}{|l|c|c|c|c|}
      \hline
      \textbf{Compiler} & \textbf{Reportet} & \textbf{Gefixt} & \textbf{Duplikate} & \textbf{Invalid} \\
      \hline
      GCC & 21 & 4 & 3 & 2 \\
      \hline
      LLVM/Clang & 26 & 7 & 0 & 1 \\
      \hline
      \textbf{Gesamt} & \textbf{47} & \textbf{11} & \textbf{3} & \textbf{3} \\
      \hline
    \end{tabular}
  \end{center}
  \vspace{0.3cm}
  \begin{itemize}
    \item Best\"atigte Bugs Gesamt: 47 (GCC + LLVM)
    \item GCC: 21 Bugs reportet \\
      4 Gefixt, 3 Duplikate, 2 Invalid, 21 Wrong Output, 0 Crashes
    \item LLVM/Clang: 26 Bugs reportet \\
      7 Gefixt, 0 Duplikate, 1 Invalid, 23 Wrong Output, 3 Crashes
    \item Status: 11 Bugs zum Zeitpunkt der Ver\"offentlichung gefixt (23\,\%)
    \item Best\"atigungsrate: 47/47 = 100\,\% (alle Bugs von Entwicklern best\"atigt)
  \end{itemize}
\end{frame}
% Meiste Bugs sind "Wrong Output"—Debugger zeigt falsche Werte.
% 3 LLDB Crashes durch malformed DWARF-Parsing.

\begin{frame}{Betroffene Compiler-Versionen}
  \begin{itemize}
    \item Betroffene GCC-Versionen:
      \begin{itemize}
        \item GCC 6.X, 7.X, 8.X, 9.X, Trunk
        \item Einige Bugs datieren zur\"uck zu GCC 4.4 (9 Jahre alt!)
      \end{itemize}
    \item Betroffene LLVM-Versionen:
      \begin{itemize}
        \item LLVM 5.X, 6.X, 7.X, 8.X, 9.X, Trunk
      \end{itemize}
    \item Implikation: Dies ist ein anhaltendes Problem, nicht historisch
    \item Latente Bugs: Einige Probleme existierten Jahre vor Entdeckung
  \end{itemize}
\end{frame}
% Bugs spannen mehrere Major-Releases.
% Einige sind latent über Jahre vor Entdeckung.
% Sogar Trunk (Development)-Versionen haben Bugs.
% Dies zeigt Bedarf für kontinuierliches Testing.
% Das Problem ist nicht gelöst—laufende Forschung nötig.

\begin{frame}{Betroffene Optimierungs-Level}
  \begin{itemize}
    \item GCC: Bugs bei \texttt{-O1}, \texttt{-O2}, \texttt{-O3}, \texttt{-Og}
      \begin{itemize}
        \item Sogar \texttt{-Og} (optimiert f\"ur Debugging) hat Bugs!
      \end{itemize}
    \item LLVM: Bugs bei \texttt{-O1}, \texttt{-O2}, \texttt{-O3}
      \begin{itemize}
        \item \texttt{-Og} ist Wrapper f\"ur \texttt{-O1} in LLVM
      \end{itemize}
    \item Wichtige Erkenntnis: Bugs existieren auf ALLEN Optimierungs-Leveln
    \item Ursache: Debug-Werte falsch von Passes aktualisiert, folgende Passes l\"oschen falsche Info nicht
    \item Implikation: Kein Optimierungs-Level f\"ur korrekte Debug-Info vertrauensw\"urdig
  \end{itemize}
\end{frame}
% Sogar debug-freundliche Optimierungs-Flags haben Probleme.
% Das Problem kaskadiert—ein Pass erstellt falsche Info, spätere Passes fixen nicht.
% Deshalb ist umfassendes Testing über alle Level wichtig.
% Entwickler können sich auf kein Optimierungs-Level für Debug-Info verlassen.

\begin{frame}{LLVM-Komponenten mit Bugs}
  \begin{columns}
    \column{0.55\textwidth}
      \begin{tabular}{|l|c|}
        \hline
        \textbf{Komponente} & \textbf{Bug-Anzahl} \\
        \hline
        DWARF Parser & 6 \\
        \hline
        SelectionDAG & 5 \\
        \hline
        SROA & 2 \\
        \hline
        InstCombine & 2 \\
        \hline
        EarlyCSE & 2 \\
        \hline
        MachineScheduler & 2 \\
        \hline
      \end{tabular}
    \column{0.45\textwidth}
      \begin{tabular}{|l|c|}
        \hline
        \textbf{Komponente} & \textbf{Bug-Anzahl} \\
        \hline
        LoopDeletion & 1 \\
        \hline
        SimplifyCFG & 1 \\
        \hline
        LSR & 1 \\
        \hline
        DeadStoreElimination & 1 \\
        \hline
        Inliner & 1 \\
        \hline
        CodegenPrepare & 1 \\
        \hline
      \end{tabular}
  \end{columns}
  \vspace{0.3cm}
  \begin{itemize}
    \item Wichtige Erkenntnis: DWARF Parser hat meiste Bugs (inklusive 3 LLDB Crashes)
    \item Verteilung: Bugs \"uber viele Optimierungs-Passes verteilt
  \end{itemize}
\end{frame}

\section{Fallstudien}


\begin{frame}[fragile]{Fallstudie: GCC wrong value bug 92387}
  \begin{lstlisting}[basicstyle=\ttfamily\tiny]
static int a[6];
int b;
static void c() { b = 5 ^ a[b & 5]; }
int main() {
    int i, d = 0;
    c();
    i = 0;
    for (; i < 8; i++) {
        c();
        if(d) printf("%d",i); // i
    }
}
  \end{lstlisting}
  \begin{itemize}
    \item Korrektheitsverletzung: $\mathit{PRINT}(P,\text{Zeile }10,i) = 0 \not\equiv 1 = \mathit{DBG}(P', A_{\langle 10,i\rangle})$ \\
    $\alert{BUG}$ (Definition verletzt)
    \item Optimierung: \texttt{-O1}; Status: Gefixt
  \end{itemize}
  \begin{center}
    \href{https://gcc.gnu.org/bugzilla/show_bug.cgi?id=92387}{\texttt{gcc.gnu.org/bugzilla/show\_bug.cgi?id=92387}}
  \end{center}
\end{frame}
% https://gcc.gnu.org/bugzilla/show_bug.cgi?id=92387

\begin{frame}[fragile]{Fallstudie: LLVM Wrong Value Bug 43957}
  \begin{lstlisting}[basicstyle=\ttfamily\tiny]
char a, b;
int main() {
    --b;
    unsigned l_801 = 36901;
    ++l_801;
    if (a) // l_801
        b = 0;
}
  \end{lstlisting}
  \vspace{0.3cm}
  \begin{itemize}
    \item Korrektheitsverletzung: $\mathit{PRINT}(P,\text{Zeile }6,l\_801) = 36902 \not\equiv 36901 = \mathit{DBG}(P', A_{\langle 6,l\_801\rangle})$ \\
    $\alert{BUG}$ (Definition verletzt)
    \item Optimierung: \texttt{-O3}; Status: Gefixt
  \end{itemize}
  \begin{center}
    \href{https://bugs.llvm.org/show_bug.cgi?id=43957}{\texttt{bugs.llvm.org/show\_bug.cgi/?id=43957}}
  \end{center}
\end{frame}
% https://bugs.llvm.org/show_bug.cgi?id=43957


\section{Fazit \& Zuk\"unftige Arbeit}


\begin{frame}[fragile]{Framework-Generalisierbarkeit -- Rust Compiler}
  \begin{lstlisting}[basicstyle=\ttfamily\tiny]
fn main() {
    let mut counter = 0;
    {
        let mut tick = || counter += 1;
        tick();
        tick();
        lib::opt_me_not();  // counter sollte 2 sein
    }
    assert_eq!(counter, 2);
}
  \end{lstlisting}
  \begin{itemize}
    \item 2 best\"atigte Bugs in 3 Tagen gefunden
    \item Bug: Rust Debugger zeigt \texttt{counter = 0}, sollte 2 sein
    \item Optimierung: \texttt{-C opt-level=1}
    \item Anpassungen n\"otig:
      \begin{itemize}
        \item Grammar-Instrumentierung (Rust vs. C)
        \item Kompilierungseinheit = crate (nicht Datei)
        \item Rust's safe Code hilft bei UB-Erkennung
      \end{itemize}
  \end{itemize}
\end{frame}
% Rust-Bugs beweisen: Framework generalisiert über C hinaus.

\begin{frame}{Fazit}
  \begin{itemize}
    \item Problem: Debug-Informationen f\"ur optimierten Code sind oft inkorrekt
    \item L\"osung: Actionable Program Framework mit formalen Definitionen
      \begin{itemize}
        \item Algorithmus: Systematisches Testing
      \end{itemize}
    \item Auswirkung: 47 Bugs in GCC/LLVM, 2 in Rust
    \item Community-Reaktion: 11 Bugs gefixt, 100\,\% Best\"atigungsrate
    \item Zuk\"unftige Arbeit:
      \begin{itemize}
        \item Debugger-Korrektheit direkt testen
        \item Auf mehr Sprachen erweitern
        \item Vollst\"andigkeit testen (fehlende Debug-Info vs. falsche Info)
      \end{itemize}
  \end{itemize}
\end{frame}


\begin{frame}{Abschlie\ss{}ende Gedanken}
  \begin{block}{LLVM-Entwickler (Bug 39896)}
    ``Thanks for the PR, this was likely a serious quality-of-debug-experience \\
    issue for anyone hit by it.''
  \end{block}
  \vspace{0.3cm}
  \begin{block}{DWARF-Komiteemitglied}
    ``[Your bug reports] seem clear and reasonably complete. Generating correct \\
    debug info for optimized code is important.''
  \end{block}
  \vspace{0.3cm}
  \begin{center}
    \href{https://www.cc.gatech.edu/~qrzhang/projects/debug/debug.html}{\texttt{cc.gatech.edu/$\sim$qrzhang/projects/debug/}}
  \end{center}
\end{frame}

\begin{frame}
  \begin{center}
    \LARGE{Vielen Dank f\"ur Eure Aufmerksamkeit!}
  \end{center}
\end{frame}


\end{document}
